\documentclass{article}
\usepackage{natbib,graphics,amsmath,kmq2,ifthen,amssymb,array,
  indentfirst}
\usepackage{rotating}
 
\renewcommand{\a}{\boldsymbol{\alpha}}
\renewcommand{\b}{\boldsymbol{\beta}}
\newcommand{\e}{\boldsymbol{\eta}}
\newcommand{\g}{\boldsymbol{\gamma}}
\newcommand{\m}{\boldsymbol{\mu}}
\newcommand{\mm}{\mathbf{m}}
\newcommand{\p}{\boldsymbol{\psi}}
\newcommand{\ph}{\boldsymbol{\phi}}
\renewcommand{\S}{\boldsymbol{\Sigma}}
\newcommand{\T}{\boldsymbol{\Theta}}
\renewcommand{\t}{\boldsymbol{\theta}}
\renewcommand{\L}{\boldsymbol{\Lambda}}
\renewcommand{\l}{\boldsymbol{\lambda}}
\newcommand{\x}{\mathbf{x}}
\newcommand{\X}{\mathbf{X}}
\newcommand{\y}{\mathbf{y}}
\newcommand{\w}{\mathbf{w}}
\newcommand{\W}{\mathbf{W}}
\renewcommand{\u}{\mathbf{u}}
\newcommand{\U}{\mathbf{U}}
\newcommand{\z}{\mathbf{z}}
\newcommand{\Z}{\mathbf{Z}}
\newcommand{\V}{\mathbf{V}}
\newcommand{\E}{\mathbf{E}}
\newcommand{\Y}{\mathbf{Y}}
\newcommand{\F}{\mathbf{F}}
\renewcommand{\I}{\mathbb{I}}

\newtheorem{property}{Property}
%\newtheorem{assumption}{Assumption}
\newtheorem{assumption}{Effect Maximizing Assumption}




%\newcommand{\ob}[1]{\textcolor[rgb]{0.65, 0.65, 0.65}{{#1}}}
\newcommand{\ob}[1]{\textcolor[rgb]{1, 0, 0}{{#1}}}


%% === hyperref options ===
\usepackage{color}
\usepackage[pdftex, bookmarksopen=true, 
bookmarksnumbered=true, linkcolor=webred]{hyperref}
\definecolor{webgreen}{rgb}{0, 0.5, 0}
\definecolor{webblue}{rgb}{0, 0, 0.5}
\definecolor{webred}{rgb}{0.5, 0, 0}


\newcommand{\indep}{{\bot\negthickspace\negthickspace\bot}}
\newcommand{\notindep}{{\bot\negthickspace\negthickspace\bot}{\negthinspace
  \negmedspace \negthickspace \negthickspace \diagup
  \;}}
\newcommand{\notone}{1 \negthinspace \negthickspace \negthickspace \diagup} 
\newcommand{\notnegone}{-1 \negthickspace \negthickspace \negthickspace \diagup}
\newcommand{\notzero}{0 \negthinspace \negthickspace \negthickspace \diagup}


\title{Gov 1000/2000/2000e Causality Lecture Notes}

\author{Adam Glynn}
\date{}

\begin{document}
\maketitle

\double

\section{Introduction}

The purpose of this document is to relate the modern, counterfactual linear-regression approach to causal inference (described in Gelman/Hill (GH) Chapter 9) to the classical approach (described in nearly every econometrics textbook including Ashenfelter, Levin, and Zimermann Chapters 9-12). This will be done within the context of the papers you have read for this class. 

Note: In composing these notes, I have assumed that you have already done the required reading for this week (especially Gelman/Hill Chapter 9). I have added page numbers where I think they will be useful. Also, these notes are very rough, do not include the appropriate citations, and should not be circulated, cited, or used for any purpose other than the course.



\section{The Classical Econometric Approach to Causal Inference with Linear Regression}


\subsection{An example where the assumptions hold (mostly) for simple linear regression} 

Consider the Kuklinski et al. (1997) experiment that we have discussed in detail over the past few weeks. We will denote an indicator variable $T_i$ which takes the value 1 if individual $i$ received the treatment list and 0 if individual $i$ received the baseline list. The outcome variable is the response to the HOW MANY question, which we will denote with variable $Y_i$. Assumption (A1) on pg. 125 of ALZ states that the outcome can be written as a linear function of the treatment, 
\begin{align*}
Y_i = \beta_0 + \beta_1 T_i + \epsilon_i
\end{align*}
and (A2) states that the average of the errors is zero ($E[\epsilon_i]=0$). As stated on ALZ pg. 127, these imply that the  expectation of the outcome is linear in the treatment, 
\begin{align*}
E[Y_i] = \beta_0 + \beta_1 T_i
\end{align*}
Notice that this assumption of the linear expectation holds trivially for this example because $T$ only takes on two values. However, we have seen examples where this assumption is not satisfied. Recall that the regression of percent liberal support in judicial power cases on the Segal/Cover scores (see Figure \ref{f:nonlinear}) seems to show a non-linear relationship. 

Assumption (A3) on ALZ pg. 127 states that the variance of the error terms is constant ($V[\epsilon_i]=\sigma^2_{\epsilon}$ for all $i$). As we discussed in lecture, this assumption probably does not hold for this example. The treatment list is the baseline list plus an additional item, therefore the variance for individuals in the treatment group should be larger than the variance for individuals in the baseline group. For example, if the first individual is in the treatment group and the second individual is in the control group, then $V[\epsilon_1] > V[\epsilon_2]$, and therefore (A3) does not hold.  However, as we have discussed, due to the sample size, we were able to relax this assumption (allowing a different variance for each group), and the standard errors do not change that much. 

Assumption (A4) on ALZ pg. 127 states that the covariance between the errors of any two individuals is zero ($Cov[\epsilon_i,\epsilon_j]=0$ for all $i \ne j$). For this example, this assumption is justified because the individuals in the Kuklinski study were randomly sampled from a large population and were therefore independent.  Furthermore, independence implies zero correlation ($\epsilon_i \bot \epsilon_j \Rightarrow Cov[\epsilon_i,\epsilon_j]=0$ for all $i\ne j$). However, in circumstances where some of the observations are repeated on the same units, this assumption will not hold. For example, suppose we use the GDP of the US between 1950 and 2000 as our outcome variable. Clearly, $GDP_{1950}$ is related to $GDP_{1951}$, and therefore the assumption will not hold.\footnote{To some extent, this can be ameliorated by including many explanatory variables, but regardless of which explanatory variables are included in the analysis, it will often be hard to argue that one has completely broken the dependence across time.}

Assumption (A5) on ALZ pg. 128 states that the covariance between the error term and the treatment is zero for each individual ($Cov[T_i,\epsilon_i]=0$ for all $i$). This is the key causal assumption, and it is justified by the randomization of assignment to the treatment and baseline lists. Random assignment implies that knowing whether someone is in the treatment or the baseline group doesn't tell us anything about their pre-treatment variables. Therefore $T_i$ is independent of $\epsilon_i$, which implies zero covariance ($T_i \bot \epsilon_i \Rightarrow Cov[T_i,\epsilon_i]=0$ for all $i$).  

Assumption (A6) on ALZ pg. 128 states that the error terms are normally distributed ($\epsilon_i \sim N(0,\sigma^2_{\epsilon})$). As we have discussed, this is unlikely to hold for this example. The responses to the lists only take on distinct values (4 for the baseline group and 5 for the treatment group). Even after we subtract $\beta_0 + \beta_1T_i$ from these values, the errors can only take on a distinct number of values, and therefore, they are not normally distributed. However, the sample size in this example is large enough that the Central Limit Theorem kicks in, so the violation of this assumption does not cause a problem. 

\subsection{An example where the assumptions hold (mostly) for multiple regression} 

Consider the Washington (2008) paper that you read this week. A simplification of the model the author presents in the first equation can be written as the following:
\begin{align*}
Y_i = \beta_0 + \beta_1 Girls_i + \beta_2 Children_i + \epsilon_i
\end{align*}

In a couple of weeks, we will discuss diagnosing violations of assumptions (A1-A4 and A6) in the multiple regression model (although simple logic would imply that we at least need an interaction in this model). For now, we will focus on (A5). When we are interested in estimating $\beta_2$ in the multiple regression setting, the key assumption is the following:
\begin{align*}
Cov[Girls_i,\epsilon_i]=0 \text{ for all } i
\end{align*}
As argued in the paper, we have some reason to believe that conditional on the number of children an individual has, the number of girls they have is essentially random.\footnote{Washington investigates a number of ways that this assumption may not hold.} If this is true, then assumption (A5) will hold for this model. 


\subsection{An example demonstrating omitted variable bias}

Suppose we believe that the two explanatory variable model is the true model,
\begin{align*}
Y_i &= \beta_0 + \beta_1 Girls_i + \beta_2 Children_i + \epsilon_i
\end{align*}
but suppose that we omit the $Children$ variable from the analysis
 \begin{align*}
Y_i &= \beta_0^{*} + \beta_1^{*} Girls_i + \epsilon_i^{*}
\end{align*}

Certainly, we would not now expect assumption (A5) to hold, because the number of girls should covary with the number of children, and the number of children is now in the error term. We can make this more explicit by writing down a model in the manner of GH pg. 170 (this is essentially equivalent to the presentation in most econometrics textbooks).  If we write the model for children as the following: 
\begin{align}\label{eq:children}
Children_i &= \gamma_0 + \gamma_1 Girls_i + \nu_i
\end{align}
then the two explanatory variable model can be re-written as the following:
\begin{align*}
Y_i &= \beta_0 + \beta_1 Girls_i + \beta_2 Children_i + \epsilon_i \\
&= \beta_0 + \beta_1 Girls_i + \beta_2 (\gamma_0 + \gamma_1 Girls_i + \nu_i) + \epsilon_i \\
&= \beta_0 + \beta_2 \gamma_0 + (\beta_1 + \beta_2 \gamma_1) Girls_i + \epsilon_i + \beta_2 \nu_i
\end{align*}
this implies that $\beta_1^{*} = (\beta_1 + \beta_2 \gamma_1)$ and $\beta_2 \gamma_1$ describes the bias due to omitting the $Children$ variable from the model. Using the 105th congress and the AAUW scores, we can see the results of this bias. In the multiple regression model, the effect of an extra girl (conditional on number of children) is estimated to increase AAUW score by almost 6 points ($\widehat{\beta}_1= 5.776$). However, if we fail to condition on the number of children, the effect of an extra girl is estimated to decrease AAUW score by almost 3 points ($\widehat{\beta}_1^{*}= -2.784$).

\subsection{A number of things that are confusing about the classical approach}

While the model(s) represented by the classical approach can be quite useful in a number of circumstances, there are certain aspects of the approach that can be confusing. Most of these stem from a lack of specificity about the meaning of the parameters (the $\beta$s and the $\gamma$s).\footnote{The following issues are not exhaustive.} 

\subsubsection{Is $\beta_1^{*}$ meant to be causal?}

As we have discussed previously, it is not clear from the models whether we are attempting to interpret $\widehat{\beta}_1^{*}$ predictively or causally. If we just want to know the expected difference in AAUW score between two representatives that happen to have difference of one in the number of girls, then -2.784 is a perfectly reasonable estimate. It is only when we want to interpret $\widehat{\beta}_1^{*}$ causally that we get into trouble.

\subsubsection{Is $\gamma_1$ meant to be causal?}

Even if we are interpreting $\beta_1$ causally, it is difficult to see how we could interpret $\gamma_1$ causally. Under the assumptions of Washington (2008), $Children$ has an effect on $Girls$, not the other way around. Therefore the model for children we used in our omitted variable bias calculation \eqref{eq:children} has the causal relationship reversed. Furthermore, the importance of this distinction is not confined to this example. \emph{If} the direction of causality really did run in the direction implied by the model in \eqref{eq:children}, we would almost certainly \emph{not} want to include the ``omitted'' variable in the model. 

To see this, consider the other paper we read this week. In Gerber, Green, and Larimer (2008), the authors used a direct mail campaign in order to assess the effects of social pressure on voter turnout. In order to simplify the discussion, we will only consider the most extreme social pressure mailing (Neighbors Treatment) and compare this treatment to the control group. In the following, we will denote an indicator variable $M_i$ which takes the value 1 if individual $i$ received the mailing and 0 if individual $i$ did not receive the mailing. We will denote the outcome variable $V_i$ which takes the value 1 if individual $i$ voted and 0 if individual $i$ did not vote.\footnote{The unit of analysis was actually households and not individuals, but we will simplify the discussion by assuming that individuals were the unit of analysis.}

Suppose that after the mailings had been sent out, the authors interviewed the subjects and asked them whether they intended to vote. We could denote this variable as $I_i$ which takes the value 1 if individual $i$ intends to vote and 0 if individual $i$ does not intend to vote. Here is the key question? Should we include the $I_i$ variable in the regression model? It is certainly correlated with the mailing treatment. And if we leave it out, it would seem to be included in the error term. Therefore, the classical econometric assumptions justify the inclusion of the variable in the model. 

However, if we want to estimate the total effect of the mailing on turnout, we should not include this variable in the model. Intuitively, including this variable will block most of the effect of the mailing, and we would get a biased estimate. This type of bias is known as post-treatment bias (see GH pg. 188). 

Figure \ref{f:post} demonstrates the logic behind omitted variable bias and post-treatment bias. In the left-hand figure, we see that the $Children$ points into both $Girls$ and $AAUWscore$. Therefore, we need to control for the $Children$ variable, blocking the path of spurious correlation. In the right-hand figure, $Post-Mailing Turnout Intention$ is on the path from $Mailing$ to $Vote$. Therefore, if we condition on $Post-Mailing Turnout Intention$, we will block the indirect path from $Mailing$ to $Vote$, thus inducing post-treatment bias. 

\subsubsection{What do we really mean by assumption (A1)?}

In fact, if we consider the social pressure example more closely, many of the classical assumptions do not appear to make sense. For example, when we consider assumption (A1) by itself, the social pressure experiment would seem to not fit this model:
\begin{align*}
V_i = \beta_0 + \beta_1 M_i + \epsilon_i
\end{align*}
Notice that $V_i$ can only take the values 0 or 1 (it is binary). Therefore, it seems impossible for this linear model to fit. However, intuitively, we know that the correct way to analyze this experimental data is to take the difference between the proportion of turnout among those that received the mailing and the proportion of turnout among those that did not receive the mailing. Because proportions are means for binary variables, this is equivalent to a linear regression of $V_i$ on $M_i$. In other words, we know that OLS is a good way to analyze the data from the social pressure experiment, so it is a bit disconcerting that this data clearly violates (A1) from ALZ.

In fact, if we interpret the $\beta_1$ in this model as an individual-level causal parameter, then we would be hard pressed to believe assumption (A1) for almost any social science application. Taken as an individual level parameter, assumption (A1) implies that the effects of the treatment are the same for all units in the analysis (e.g., it assumes that the mailing affects literate and illiterate individuals in the same manner). This constant-effects assumption is not necessary, and this is why I have mostly jumped straight to the linear expectation assumption in the lecture notes: 
\begin{align}\label{eq:A1plusA2}
E[V_i] = \beta_0 + \beta_1 M_i
\end{align}
Notice that while the ALZ assumption (A1) is clearly violated by the data in the social pressure experiment, the linear CEF of \eqref{eq:A1plusA2} is trivially satisfied by the data in in the social pressure experiment. 

As a side note, assumption (A1) is the source of much confusion. As we will discuss in detail in the next section, we do not need to assume constant effects in order for linear regression to provide unbiased estimates of average causal effects. However, without the constant-effects assumption, many more advanced techniques (e.g., instrumental variables) do not provide even consistent estimates of average causal effects.\footnote{With certain assumptions, instrumental variables techniques sometimes provide consistent estimates for a weighted averages of effects. Whether these weighted averages are of theoretical interest is the subject of much debate.} Furthermore, it is important to note that uncertainty over whether or not we are making the constant-effects assumption has led to much of the confusion between the quantitative and qualitative communities. 


\subsubsection{What happens if we include an irrelevant explanatory variable in the regression?}

Many econometrics textbooks claim that including an irrelevant explanatory variable in the regression cannot induce bias. It turns out that this claim depends subtly on exactly what one means by irrelevant. The challenge problem on the problem set will explore this. 





\begin{figure}[!t]
  \begin{center}
   \includegraphics[width=6.5in, angle=0]{residuals3.pdf}
  \end{center}
  \caption{\emph{An example demonstrating a violation of the linearity assumption (A1,A2 from ALZ).} The correct model here appears to be curved and not linear.}\label{f:nonlinear}
\end{figure}

\begin{figure}[!t]
  \begin{center}
   \includegraphics[width=6.5in, angle=0]{GirlsVotes.pdf}
  \end{center}
  \caption{\emph{A graphical illustration of omitted variable bias and post-treatment bias.} If the model on the left is the true model and we omit $Children$ from our regression, we will get omitted variable bias in our estimate of the $Girls$ effect. If the model on the right is the true model and we include $Post-Mailing$ $Turnout$ $Intention$ in the regression, we will get post-treatment bias in our estimate of the $Mailing$ effect.}\label{f:post}
\end{figure}

\clearpage
\newpage


\section{The Counterfactual Approach to Causal Inference with Linear Regression}


\begin{quote}
{\it In the beginning, we should ask, What is the causal relationship of interest?}
\end{quote}
\begin{center}
Angrist and Pischke (2009), pg. 3
\end{center}


\subsection{The Fundamental Problem of Causal Inference}

The modern approach to causal inference frames causal relationships in terms of observed outcomes and counterfactual outcomes. The fundamental problem of causal inference is that we cannot observe the counterfactual outcomes, however, the difficulty of this problem varies depending on the variable types. 


\subsubsection{Dichotomous Treatments and Dichotomous Outcomes}

The fundamental problem of causal inference is easiest to solve when we have a single dichotomous treatment and a single dichotomous outcome. For example, in the social pressure experiment of Gerber, Green, and Larimer (2008), many individuals receive a mailing that socially pressures them to vote.\footnote{In order to simplify the discussion, we will ignore the other types of mailings, and we will refer to households as individuals.} It is conventional to label this active treatment condition as \emph{treatment} (which we code as  $T=1$). Other individuals in the study did not receive a mailing. It is conventional to label this passive treatment condition as \emph{control} (which we code as $T=0$).\footnote{This multiple use of the word treatment can be confusing, but the meaning is usually clear from context, and it has become standard usage within the field of causal inference.}  The social pressure experiment was designed to investigate the causal effects of this mailing on turnout, and therefore the outcome variable was whether or not each individual voted. 


As in Gelman/Hill pg. 178, if we are willing to assume SUTVA (Stable Unit Treatment Value Assumption), that an individual's outcome (e.g. decision to turnout to vote) does not depend on whether other individuals receive the treatment or control condition, and we further assume that each individual could potentially receive either the treatment or the control condition (i.e., receive or not receive the mailing), then we can categorize the individuals in our study into one of four types: \emph{Always} voters would have voted regardless of whether they received the treatment or control condition (e.g. mailing or no mailing). \emph{Never} voters would have failed to vote regardless of whether they received the treatment or control condition. \emph{Caused} individuals would have voted if they received the treatment condition (mailing) but would have failed to vote if they had received the control condition (no mailing). \emph{Prevented} individuals would have failed to vote if they received the treatment condition (mailing) but would have voted if they had received the control condition (no mailing). 

These response types are summarized in Table \ref{t:types}, where we have labeled voting outcomes with a \emph{1} and non-voting outcomes with a \emph{0}. Furthermore, this quantification of outcomes allows the definition of individual treatment effects (ITEs) as the difference between the potential outcome under treatment (denoted with $y^1$) and the potential outcome under control (denoted with $y^0$). This is the definition provided on GH pg. 170, although we have used a dichotomous outcome for this example. 

%Note that we might say that treatment is a sufficient condition for an always individual and an insufficient condition for a never individual, and we could also say that treatment is a necessary and sufficient condition for a caused individual and an unnecessary and insufficient condition for a prevented individual. However, once we observe the outcome under one of the treatment conditions, we will find that distinctions may not conform directly with the typical usage of these terms. 

\begin{table}
\begin{center}
\caption{\label{t:types}Potential outcomes, response types, and individual treatment effects (ITEs) with a dichotomous treatment and dichotomous outcome.}
\begin{tabular}{rr | r | r}
\hline \hline
\multicolumn{2}{c|}{Potential Outcomes} & & \multicolumn{1}{c}{ITE}\\
$y^1$ & $y^0$ & Response Type & $y^1 - y^0$ \\
\hline
voted (1) & voted (1) & always & 1-1=0 \\
didn't vote (0) & didn't vote (0) & never & 0-0=0 \\
voted (1)& didn't vote (0) & caused & 1-0=1  \\
didn't vote (0) & voted (1)& prevented & 0-1=-1 \\
\hline \hline
\end{tabular}
\end{center}
\end{table}


\begin{table}
\begin{center}
\caption{\label{t:ideal} The data if potential outcomes were known for all units in our study.}
\begin{tabular}{r|rrr|r|r|r|r}
$Unit_i$ & $X_{i1}$ & $X_{i2}$ & $X_{i3}$ & $T_i $ & $Y_i^0$ & $Y_i^1$ & $Y_i^1-Y_i^0$ \\
\hline
1 & 2 & 1 & 50 & 0 & 0 & 0 & 0 \\ 
2 & 3 & 1 & 98 & 0 & 0 & 1 & 1 \\ 
3 & 2 & 2 & 80 & 1 & 1 & 0 & -1 \\ 
4 & 3 & 1 & 98 & 1 & 0 & 1 & 1\\
\vdots &  \vdots &\vdots &\vdots &\vdots &\vdots &\vdots & \vdots  \\
\end{tabular}
\end{center}
\end{table}





In an ideal research world, we would observe these types for each individual in the study (see Table \ref{t:ideal}). Unfortunately, we will never \emph{directly} observe the response type of any individual because we only observe each individual under either the treatment condition or the control condition. However, we can always rule out two of the types. If an individual didn't receive the mailing and didn't vote, then we know that their type is never or caused. If an individual didn't receive the mailing and voted, then we know that their type is always or prevented. If an individual received the mailing and didn't vote, then we know that their type is never or prevented. If an individual received the mailing and voted, then we know that their type is always or caused. Table \ref{t:anhh} summarizes these conditions as well as the associated possible values for the individual treatment effect. Table \ref{t:real} presents the data we actually have.

\begin{table}
\begin{center}
\caption{\label{t:anhh} Possible response types and individual treatment effects (ITEs) conditional on the observed treatment condition and outcome.}
\begin{tabular}{rr | r | r}
\hline \hline
\multicolumn{2}{c|}{Observed} & & \\
Treatment Condition ($T$) & Outcome ($Y$) & Possible Response Types & Possible ITEs \\
\hline
no mailing (0) & didn't vote (0) & never or caused & \{0,1\} \\
no mailing (0)& voted (1) & prevented or always & \{-1,0\} \\
mailing (1)& didn't vote (0) & prevented or never & \{-1,0\} \\
mailing (1)& voted (1) & always or caused & \{0,1\} \\
\hline \hline
\end{tabular}
\end{center}
\end{table}


\begin{table}
\begin{center}
\caption{\label{t:real} The data we actually observe.}
\begin{tabular}{r|rrr|r|r|r|r}
$Unit_i$ & $X_{i1}$ & $X_{i2}$ & $X_{i3}$ & $T_i $ & $Y_i^0$ & $Y_i^1$ & $Y_i^1-Y_i^0$ \\
\hline
1 & 2 & 1 & 50 & 0 & 0 & ? & ? \\ 
2 & 3 & 1 & 98 & 0 & 0 & ? & ? \\ 
3 & 2 & 2 & 80 & 1 & ? & 0 & ? \\ 
4 & 3 & 1 & 98 & 1 & ? & 1 & ?\\
\vdots &  \vdots &\vdots &\vdots &\vdots &\vdots &\vdots & \vdots  \\
\end{tabular}
\end{center}
\end{table}


%some of the necessary and sufficient conditions that were discussed above do not seem to conform with typical usage. For example, if we observe a success under the treatment condition, then we note that according to the previous rules, the treatment must be deemed sufficient (because treatment was deemed sufficient for both always and caused individuals). However, ... 


\subsubsection{Relationship to small-n inference and case studies}

Although it is not possible to directly observe the response type or ITE for any individual, in some cases it may be possible to collect additional information on each individual in order to guess their response type. There are at least three different approaches to this.

First, for individuals that were observed in two elections (one under the treatment condition and one under the control condition) we could compare their voting behavior across the two elections. That is, if we observe someone who did not vote in the previous election, and did not vote in the current election (after receiving the mailing), then we might assume that the mailing had no effect on them (never voter instead of prevented). Of course this conclusion assumes that this voter did not change what their behavior would have been without the mailing. More formally, this sort of crossover study assumes that each individual did not change their response type or ITE across the two elections.\footnote{In a large-$n$ study, we might relax this assumption in a number of ways. We will discuss this during the last lecture.} 

Second, we could collect additional information about individuals to examine the process by which the treatment would potentially affect the outcome. For example, suppose that we interviewed an individual that received the mailing and voted, but they claim to recycle anything that looks like junk mail, and they do not recall reading the mailing. For this individual, we might conclude that they were an \emph{always} voter. This type of analysis sometimes falls under the rubric of process tracing. 

Finally, if we find two voters (one that received the mailing and one that didn't) and we believe them to be causally homogenous with respect to the social pressure treatment (i.e., they have the same response type in Table \ref{t:anhh}), then we can determine the response type for each individual. For example, if the individual that received the mailing voted ($y^1 = 1$) and the individual that didn't receive the mailing didn't vote ($y^0 = 0$), and we believe that these two individuals have the same potential outcomes and response types, then we have determined that both individuals have an ITE of $y^1 - y^0 = 1$ (that they were caused to vote by the treatment). This is essentially Mill's method of difference. 

Notice that all of these techniques for establishing ITEs are more likely to be applied in small or medium $n$ settings. 


\subsubsection{Dichotomous Treatments and Continuous Outcomes}

With dichotomous treatments and dichotomous outcomes, the counterfactual potential outcome can only take on two possible values. For example if $T=0$, then $y^1$ is counterfactual, but we know that it can only take on two possible values $\{0,1\}$. With continuous outcomes, the number of possible values for the counterfactual potential outcome grows, and therefore the fundamental problem of causal inference becomes more difficult.

For example, in the Fish (2002) paper, the author examines the country-level effects of having an Islamic religious tradition ($T=1$) versus not having an Islamic religious tradition ($T=0$) on the level of democracy for a country. Level of democracy is measured in two ways (reversed Freedom House score and polity score), but both of these measures are continuous.\footnote{At least they are conceptually continuous, although in practice they can only take on a finite number of values.} Using the reversed Freedom House score as our outcome variable, Table \ref{t:DiCon} presents the first five entries from the data set. The first country (Afghanistan) has an Islamic religious tradition ($T=1$) and  a democracy score of $Y=1.10$. Therefore,  $y^1= 1.10$ for this country, and $y^0$ could be anything within the [1 --- 7] range. This implies that the ITE can be anything within the [-5.90 --- 0.10] range. The other observations in the data set can be analyzed similarly. 

There are two important things to note about the continuous outcomes (as opposed to dichotomous outcomes). First, some of the small-$n$ techniques become more difficult to implement. For example, if we attempt to use process tracing to establish the democracy score that Afghanistan would have had, had it not had an Islamic religious tradition, we must now choose a number on a seven point scale (as opposed to the social pressure example, where we only had to determine whether a voter would not have voted). This is one of the reasons that small-$n$ studies tend to dichotomize (or sometimes trichotomize) variables that are conceptually continuous. Second, if we are willing to resign ourselves to estimating average effects, continuous outcomes may actually make large-$n$ analysis easier (because the linear expectation assumption will be more likely to hold when we have continuous explanatory variables).  

\begin{table}
\begin{center}
\caption{\label{t:DiCon} Possible individual treatment effects (ITEs) conditional on the observed treatment condition and outcome when the outcome is a continuous variable that ranges between 1 and 7 (as it does for the reversed Freedom House score used in Fish (2002)).}
\begin{tabular}{rrr | rr | r}
\hline \hline
\multicolumn{3}{c|}{Observed} & & & ITE \\
Country& Islamic Tradition ($T$) & Democracy ($Y$) & $y^1$ & $y^0$ & $y^1 - y^0$  \\
\hline
Afghanistan & 1  & 1.10 & 1.10 & [1 --- 7] & [-5.90 --- 0.10] \\
Albania & 1 & 4.10 & 4.10 & [1 --- 7]  & [-2.90 --- 3.10] \\
Algeria & 1 & 2.15 & 2.15 & [1 --- 7]  & [-4.85 --- 1.15]  \\
Angola & 0  & 1.90 &  [1 --- 7]  & 1.90 & [-0.90 --- 5.10] \\
Argentina & 0 & 5.65&  [1 --- 7]  & 5.65  & [-4.65 --- 1.35] \\
$\vdots$ & $\vdots$ & $\vdots$ & $\vdots$ & $\vdots$ & \vdots  \\
\hline \hline
\end{tabular}
\end{center}
\end{table}


\subsubsection{Continuous Treatments and Multiple Treatments}

While a continuous outcome (with a dichotomous treatment) increases the number of possible values for the counterfactual potential outcome, continuous treatments increase the number of counterfactual potential outcomes. For example, we might think of income in the Fish (2002) paper as a treatment. In the Fish (2002) data, income values range between 2 and 4.662, so we could potentially be interested in many, many potential outcomes (e.g., \{$y^2$,$y^{2.001}$, ... ,$y^{4.661}$, $y^{4.662}$\}). Furthermore, for any country we only observe one of these outcomes. Afghanistan has an income of $2.25$ and an democracy score of $1.10$, so we know that $y^{2.250}=1.10$, but we don't know the value of any of the other potential outcomes in the set of \{$y^2$,$y^{2.001}$, ..., $y^{4.661}$, $y^{4.662}$\}. The fundamental problem of causal inference has become much more difficult to solve. 

There are essentially two strategies for dealing with this problem. First, we can dichotomize the treatment variable (i.e., rich countries and poor countries) or break it up into a manageable number of groups. Unfortunately, this tends to make the SUTVA assumption less plausible because different values of the continuous treatment will be considered the same once the variable is dichotomized. The second strategy is to assume some amount of smoothness between the potential outcomes. For example, with a continuous treatment and a continuous outcome, we will often assume that there is a linear relationship between the potential outcomes for any given unit (we still allow heterogeneity across units). For example, if we assume linearity, and I tell you that for Afghanistan $y_{Afg}^{3.250}=1.25$, then we would know that$y_{Afg}^{4.250}=1.40$ because $y_{Afg}^{2.250}=1.10$. That is, an increase of 1 in income leads to an increase of $0.15$ in democracy score for Afghanistan. And an increase of 1 in income might lead to a different increase in democracy score for the other countries in the data set.  


This all demonstrates that the fundamental problem of causal inference becomes quite difficult to solve with continuous (or many valued treatments), and the only way to proceed is to make strong and untestable assumptions. One sign of the fundamental difficulty with continuous treatments is the cost required to conduct an experiment that gets at the full range of a continuous variable. For example, suppose we were interested in the effects of a aspirin on blood pressure, and we wanted to get a sense of these effects over the full range from 250mg per day to 1000mg per day. The cost would be prohibitive if we tried to have 751 treatment arms in our study (one for each mg level). In practice, we are often happy with experiments with four treatment arms (250mg, 500mg, 750mg, 1000mg). 

Similarly, the fundamental problem of causal inference also becomes quite difficult to solve if we want to consider many competing explanations for a particular outcome. Suppose we have five competing explanations (this is a typical number for many areas of research), and we are willing to dichotomize the five associated variables. This corresponds to $2^5=32$ different treatment arms to our study, which would require a massive sample size in a randomized experiment (much of the experimental design literature is devoted to getting around this problem). If you want to test between competing explanations in small-$n$ non-randomized studies, you will have to make a number of strong assumptions, and rely heavily on process tracing.  \\


\single
\noindent Discussion Questions:
\begin{itemize}
\item Does the SUTVA assumption seem plausible for the Gerber, Green, and Larimer (2008) paper and/or the Fish (2002) paper?
\end{itemize}
\double

\clearpage

\subsection{The Experimental ``Solution'' to the Problem of Causal Inference}


\begin{quote}
{\it The second research FAQ is concerned with the experiment that could ideally be used to capture the causal effect of interest?}
\end{quote}
\begin{center}
Angrist and Pischke (2009), pg. 4
\end{center}


If we believed that all of the units of analysis (the individuals, representatives, or countries) were causally homogenous (i.e., had the same values for their potential outcomes), then we could solve the fundamental problem of causal inference by comparing the outcomes of units with different values of the treatment. Unfortunately, outside of the physical sciences we would almost never believe that the units were actually homogenous.\footnote{Galileo's pushing-versus-dropping experiment is a canonical example where we might believe homogeneity. As the story goes, Galileo pushed one cannonball off a tower at the same time as dropping an identical cannonball off the tower, and they hit the ground at the same time. The reason we find this experiment convincing is that we believe the cannonballs actually were identical in terms of their potential outcomes (falling times under the pushing and dropping conditions).} 

We might try to relax the homogeneity assumption by matching pairs of units (one receives the treatment and one receives the control) that appear to be identical (prior to the application of the treatment/control). If we believe that each of the matched pairs has the same values of the potential outcomes, then we could solve the fundamental problem of causal inference by comparing the outcomes of units within each matched pair. Unfortunately, no matter how many pre-treatment measurements we make on the units, we will never know if we have failed to measure an important factor.  

Random treatment assignment solves the problem of unmeasured factors, but it only solves it on average. If we randomly assign the treatment and control conditions to our units, then prior to the application of treatment, we know that, on average, the treated units were the same as the control units for both measured \emph{and} unmeasured factors. In more precise terminology, this means that the population distribution of the potential outcomes is the same for the treated and control units, and therefore assigned treatment is independent of the distribution of potential outcomes. Another way to say this is that the assigned treatment is \emph{ignorable} in terms of the potential outcomes. In mathematical notation, this ignorability condition is typically written as $\{y^0,y^1\} \bot T$ (see GH pg. 183). 

How does this help us? Recall that our goal in regression analysis is to estimate the conditional expectation function (CEF). For example when we regress an outcome on a single treatment variable, we are attempting to estimate a conditional expectation function $E[Y|T=t]$. If $T$ takes only two values, then this function can be written as a linear function without making any assumptions:
\begin{align*}
E[Y|T=t] = \beta_0 + \beta_1 \cdot t
\end{align*}
Therefore, $\beta_1$ can be interpreted, predictively, as the difference between two conditional expectations $E[Y|T=1]-E[Y|T=0]$. Randomization allows us to additionally interpret $\beta_1$ causally (in terms of potential outcomes): 
\begin{align}
\beta_1 &= E[Y|T=1]-E[Y|T=0] \notag\\
&= E[y^1|T=1]-E[y^0|T=0] \label{eq:consistency}\\
&= E[y^1]-E[y^0] \label{eq:ignore}\\
&= E[y^1-y^0] \label{eq:propExp}
\end{align}
For a randomized experiment, $\beta_1$ is the average of the individual treatment effects, known as the average treatment effect (ATE). The justification for this causal interpretation can be seen in the three labeled equalities.

The equality in \eqref{eq:consistency} holds due to SUTVA. SUTVA allows us to state that the observed $Y$ is equal to the potential outcome under treatment ($y^1$) for those that receive the treatment ($T=1$) and is equal to the potential outcome under control ($y^0$) for those that receive the control ($T=0$). This is sometimes summarized with the equation $Y =  y^1 \cdot T + y^0 \cdot (1-T)$.
 
The equality in \eqref{eq:ignore} states that the average $y^1$ from the subpopulation that receives treatment is equal to the average $y^1$ in the overall population, and similarly that the average $y^0$ from the subpopulation that receives control is equal to the average $y^0$ in the overall population. This equality holds due to ignorability, which implies that $T$ is independent of the potential outcomes (recall that when two random variables, $V$ and $W$ are independent, $E[V|W=w] = E[V]$ for all possible $w$). 

%It is also possible to get a more intuitive sense of \eqref{eq:ignore} for large populations. If we randomly assign half of the population to receive $T=1$, then the subpopulation with $T=1$ is just a random sample from the population, and if half of the population is large (think $\frac{N}{2} \rightarrow \infty$), we would expect the average $y^1$ from this subpopulation to equal the average $y^1$ in the overall population.\footnote{It is the consistency of the sample average that allows us to reach this conclusions.} Similarly, if we randomly assign half of the population to receive $T=0$, then the subpopulation with $T=0$ is just a random sample from the population, and if half of the population is large, we would expect the average $y^0$ in this subpopulation to equal the average $y^0$ in the overall population.\footnote{For a small population, the equality in \eqref{eq:ignore} will not hold exactly under the interpretation stated above. For example, suppose our population is the 50 states in the US, the outcome variable is statewide voter turnout, the treatment condition is having election day registration (EDR) available, and control is not having EDR available. In this case, even if we could randomly assign 25 of the states to treatment and 25 of the states to control, it is unlikely that the average voter turnout among the treated states \emph{for any given randomization} would equal the average voter turnout for all 50 states if all states received the treatment. However, the equality in \eqref{eq:ignore} will hold exactly across \emph{all possible randomizations}, and it turns out that any across randomization variance can be accounted for (approximately) by positing an infinite population at the estimation stage. Specifically, if the sample/population is not too small, treating the sample as if it had been randomly sampled from a very large population provides a good approximation to what we would get if we performed the correct finite-population inference (where the randomness comes from the randomized treatment assignment). For example, even if the population of interest is the 50 states, treating those states as if they had been randomly sampled from some infinite superpopulation of states provides a good approximation of performing exact finite-population inference where we would take into account the randomness inherent in randomizing 25 of the states to receive EDR. We do not need to believe that the superpopulation of states is a real thing, it merely functions to provide an approximation.}

The equality in \eqref{eq:propExp} states that the average potential outcome under treatment minus the average potential outcome under control is equal to the average of the individual treatment effects (ATE). This equality holds because of the properties of expectations ($E[V] - E[W] = E[V-W]$).\footnote{It is important to note that this equality would not hold for other summaries of the distribution of effects. For example if we wanted to know the median effect ($median(y^1- y^0)$), it cannot be broken up in this manner. 
\begin{align*}
median(y^1- y^0) &\ne median(y^1) - median(y^0)
\end{align*}
In fact, this lack of decomposability for summaries of the distribution of effects other than the average (e.g., the median effect, the lower and upper quartile of effects, the variance of effects, etc.) has important implications for the limits of causal knowledge.}

The social pressure experiment of Gerber et al. (2008) is a good political science example of the randomized experimental approach to causal inference. The randomization of households to receive different types of mailings or no mailing and the large sample size ensure that the potential outcomes in each of the treatment groups will be the same (on average). In fact, Gerber et al. (2008) presents direct evidence of this in Table 1, showing that the average values of household size, past electoral participation, gender, and age are nearly identical in the different treatment groups. This means that if we included these variables in the regression, they would not substantially affect our estimates of the treatment effects. Furthermore, it is more important to notice that the balance on these potential control variables was achieved through random treatment assignment and not through directly matching households on these variables. Hence, we may have more confidence that the different treatment groups will be similar (on average) for any important pre-treatment factors that are not included in Table 1.\footnote{Notice as well that this table (or similar displays such as balance plots) can be produced without looking at the data on the outcome variable. That is, it is possible to debate about whether we have achieved balance on the important causal factors \emph{before} we know the results of the analysis.} This confidence allows us to use a simple regression to assess the effect of each treatment in relation to the control group:
\begin{align*}
E[Turnout] &= \beta_0 + \beta_1 \cdot Civic + \beta_2 \cdot Hawthorne + \beta_3 \cdot Self  + \beta_4 \cdot Neighbors
\end{align*}
\begin{center}
\begin{tabular}{r|r|r|r|r}
$\widehat{\beta}_0$ & $\widehat{\beta}_1$ & $\widehat{\beta}_2$ & $\widehat{\beta}_3$ & $\widehat{\beta}_4$ \\
\hline
29.7\% & 1.8\% = 31.5\% - 29.7\% & 2.5\% = 32.2\% - 29.7\% & 4.8\% = 34.5\% - 29.7\% & 8.1\% = 37.8\% - 29.7\%  \\
\hline
\end{tabular}
\end{center}
Notice that even now that we have included all of the treatments implicit in the social pressure paper, the assumptions of linearity and additivity are still trivially satisfied. As can be seen in the table, running the regression in this manner is equivalent to calculating means for the five different treatments, and then calculating the differences between the four ``mailings'' means and the control group (see ALZ pg. 177 for details).\\

\noindent Discussion Questions: 
\single
\begin{itemize}
\item What assumptions would we be making about variances if we ran this regression in R or Stata? 
\item Would those assumptions be reasonable in this case? Could we relax them?
\end{itemize}
\double

\subsection{The Multiple Regression ``Solution'' to the Problem of Causal Inference}

\begin{quote}
{\it What is your identification strategy?}
\end{quote}
\begin{center}
Angrist and Pischke (2009), pg. 7
\end{center}

If we don't have a randomized experiment, then we need to have another strategy for identification (solving the problem of causal inference). Multiple regression provides one such strategy.


As discussed in the previous section, random treatment assignment ensures that the different treament groups are comparable on both measured and unmeasured factors. Multiple regression works as a strategy of causal inference to the extent that we believe that the treatment is approximately randomized, conditional on the measured factors. In more precise terminology, this means that the population distribution of the potential outcomes is the same for the treated and control units within levels of the conditioning variables, and therefore, assigned treatment is conditionally independent of the distribution of potential outcomes. Another way to say this is that the assigned treatment is conditionally \emph{ignorable}. In mathematical notation, if $X$ is a control variable, this conditional ignorability condition is typically written as $\{y^0,y^1\} \bot T|X$ (see GH pg. 183). The Washington (2008) paper provides a good example of this strategy. Conditional on the number of children that a representative has, we might believe that the number of daughters the representative has is essentially random (due to the fact that the sex of each child is essentially random).  

In mathematical terms, this strategy works in the following manner. When we regress an outcome on a treatment variable and a conditioning variable, it is typical to assume an additive linear form:
\begin{align*}
E[Y|T=t,X=x] = \beta_0 + \beta_1 \cdot t + \beta_2 \cdot x
\end{align*}
In this model, $\beta_1$ can be interpreted predictively, as the difference between two conditional expectations $E[Y|T=t+1,X=x]-E[Y|T=t,X=x]$ for any value of $X=x$. Randomization allows us to interpret $\beta_1$ causally (in terms of potential outcomes): 
\begin{align}
\beta_1 &= E[Y|T=t+1,X=x]-E[Y|T=t,X=x] \text{ for all } t,x \label{eq:linAdd} \\
&= E_X\left\{E[Y|T=t+1,X=x]-E[Y|T=t,X=x]\right\} \text{ for all } t \label{eq:add} \\
&= E_X\left\{E[y^{t+1}|T=t+1,X=x]-E[y^{t}|T=t,X=x]\right\} \text{ for all } t \label{eq:condconsistency}\\
&= E_X\left\{E[y^{t+1}|X=x]-E[y^{t}|X=x]\right\} \text{ for all } t \label{eq:condignore}\\
&= E_X\left\{E[y^{t+1}-y^{t}|X=x]\right\} \text{ for all } t \label{eq:condprop}\\
&= E[y^{t+1}-y^{t}]  \text{ for all } t \label{eq:totalexp}
\end{align}

The equality in \eqref{eq:linAdd} holds due to the assumptions of linearity and additivity. In particular, linearity allows us to assume that moving from $T=0$ to $T=1$ has the same effect as moving from $T=1$ to $T=2$. Therefore, we can just write $T=t+1$ and $T=t$ for all $t$. 

The equality \eqref{eq:add} holds due to the assumption of additivity. The difference between two conditional expectations $E[Y|T=t+1,X=x]-E[Y|T=t,X=x] = \beta_1$ does not depend on the value of $x$. Noting that $E_X\{ \cdot \}$ is mathematical notation that means averaging over the distribution of $X$, if we average this difference ($E[Y|T=t+1,X=x]-E[Y|T=t,X=x]=\beta_1$) over the distribution of $X$, we are taking an average of a constant, which results in that constant ($\beta_1$). 

The equality in \eqref{eq:condconsistency} still holds due to SUTVA. SUTVA is a statement about individual potential outcomes, so conditioning does not really change this assumption.\footnote{However, if our goal in conditioning is to say something about joint effects, then we must make a slightly different SUTVA assumption. This will be covered in Gov 2002.} 

The equality in \eqref{eq:condignore} states that the average $y^{t+1}$ from the subpopulation that receives treatment $T=t+1$ and has $X=x$ is equal to the average $y^{t+1}$ in the subpopulation that has $X=x$, and similarly that the average $y^{t}$ from the subpopulation that receives $T=t$ and has $X=x$ is equal to the average $y^{t}$ in the subpopulation that has $X=x$. This equality now holds due to \emph{conditional} ignorability.

Finally, the equality in \eqref{eq:condprop} holds due to the properties of expectations, and the equality in \eqref{eq:totalexp} holds due to the law of total expectation. Within the Washington (2008) example for the 105th congress and the AAUW score, we estimate $\widehat{\beta}_1=5.78$. That is, the average effect of an additional daughter is estimated to be an increase of $5.78$ in AAUW score.

However, it is important to note that while the linearity and additivity assumptions were trivially satisfied for the Gerber, Green, Larimer (2008) example, these assumptions must be examined more closely for the Washington (2008) example. For now, we will examine the additivity assumption. The linearity assumption will be addressed in two weeks. 

Suppose we allow an interaction between the treatment and the conditioning variable
\begin{align*}
E[Y|T=t,X=x] = \beta_0 + \beta_1 \cdot t + \beta_2 \cdot x+ \beta_3 \cdot t \cdot x
\end{align*}
In this model, $\beta_1$ can \emph{no longer} be interpreted predictively as the difference between two conditional expectations. Instead, we predictively interpret $(\beta_1 + \beta_3 \cdot x)$ as the difference between two conditional expectations $E[Y|T=t+1,X=x]-E[Y|T=t,X=x]$ which depends on the value of $X=x$. This can be most easily seen by rearranging terms in the conditional expectation:
\begin{align*}
E[Y|T=t,X=x] = \beta_0 + \bs{(\beta_1 + \beta_3 \cdot x)} \cdot t + \beta_2 \cdot x
\end{align*}
Conditional ignorability allows us to interpret $(\beta_1 + \beta_3 \cdot x)$ causally (in terms of potential outcomes): 
\begin{align}
(\beta_1 + \beta_3 \cdot x) &= E[Y|T=t+1,X=x]-E[Y|T=t,X=x] \text{ for all } t,x \notag \\
&= E[Y|T=t+1,X=x]-E[Y|T=t,X=x] \text{ for all } t \notag \\
&= E[y^{t+1}|T=t+1,X=x]-E[y^{t}|T=t,X=x] \text{ for all } t \notag\\
&= E[y^{t+1}|X=x]-E[y^{t}|X=x] \text{ for all } t \notag \\
&= E[y^{t+1}-y^{t}|X=x]\ \text{ for all } t \label{eq:condeff}
\end{align}
After relaxing the assumption of additivity, but keeping the assumptions of linearity, SUTVA, and conditional ignorability, $(\beta_1 + \beta_3 \cdot x)$ can be interpreted as the effect of increasing $T$ by one unit, conditional on $X=x$. Within the Washington (2008) example for the 105th congress and the AAUW score, we estimate $(\widehat{\beta}_1 + \widehat{\beta}_3 \cdot x) = (10.88 - 1.41 \cdot x)$. That is, if a representative has one child, the average effect of that child being a girl is an increase of $9.47$ in AAUW score. Furthermore, the effect of an additional girl decreases as the number of children increases. For example, if a representative has two children, the average effect of changing from zero to one girl (or from one to two girls) is an increase of $8.06$ in AAUW score.

However, although the average treatment effect conditional on $X=x$ ($E[y^{t+1}-y^{t}|X=x]$) may be of substantive interest, our stated goal was to estimate the overall average treatment effect ($E[y^{t+1}-y^{t}]$). Fortunately, the law of total expectation, and our linearity assumption allow us to interpret $(\beta_1 + \beta_3 \cdot E[X])$ as the average treatment effect. 
\begin{align*}
\beta_1 + \beta_3 \cdot E_X\left\{X \right\} &= E_X\left\{\beta_1 + \beta_3 \cdot X \right\} \\
&= E_X\left\{E[Y|T=t+1,X=x]-E[Y|T=t,X=x]\right\}\\
&= E_X\left\{E[y^{t+1}|T=t+1,X=x]-E[y^{t}|T=t,X=x]\right\} \\
&= E_X\left\{E[y^{t+1}|X=x]-E[y^{t}|X=x]\right\}\\
&= E_X\left\{E[y^{t+1}-y^{t}|X=x]\right\} \\
&= E[y^{t+1}-y^{t}] 
\end{align*}
For the Washington (2008) example, this means that we can estimate the average treatment effect of an additional daughter to be $(\widehat{\beta}_1 + \widehat{\beta}_3 \cdot \overline{x}) =  (10.88 - 1.41 \cdot 2.49) = 7.37$, where $2.49$ is the average number of children in the 105th congress. Notice that by relaxing the additivity assumption, we have increased our estimate of the average effect of an extra daughter on AAUW score ($5.78$ in the additive model). Furthermore, it is easy enough to get standard errors for our new estimate. We need only re-center the number of children variable in the data set by subtracting $2.49$ from the number of children for each representative, then re-run the interaction model. The correct estimate and standard errors will now be presented for the coefficient on the number of girls. \\

\single

\noindent Discussion Questions:
\begin{itemize}
\item Notice that both of the estimates ($5.78$ and $7.37$) generated here are a good deal larger than the estimate presented in the Washington (2008) paper for the AAUW scores in the 105th congress. Why do you think this is happening?
\item What are some threats to ignorability in the Fish (2002) analysis?
\end{itemize}






%For some applications, we will only be interested combinations of the proportions of response types described in the previous section. In the program evaluation setting, we often want to know the average effect of the program on the participants. This average effect, sometimes called the Average Treatment effect on the Treated (ATT). The ATT is just the average of the ITEs for all individuals in the population that received the treatment.In the social pressure example, ATT describes the proportion caused to vote by social pressure minus the proportion caused to not vote by social pressure, and hence this is the increase in voter turnout among those that received the mailing that is due to the mailing.\footnote{With a dichotomous outcome, ATT can always be described as the difference between the proportion caused among the treated minus the proportion prevented among the treated.}  

%If instead of evaluating an existing program, we consider extending a program to the control units, then we might like to know the average effect of the program among the non-participants. This average effect is sometimes called the Average Treatment effect on the Controls (ATC). In the social pressure example, ATC describes the proportion that would have been caused to vote by the mailing minus the proportion that would have been caused to not vote by the mailing, and hence this is the increase in voter turnout among the controls that we would have seen if they had been sent the mailing.\footnote{With a dichotomous outcome, ATC can always be  can be described as the proportion caused minus the proportion prevented among the control units.} 

%Finally, if instead we want to compare the average effect of the policy that treats all units in the population (send the mailing to everyone) to the policy that treats no units in the population (send the mailing to no one), then we need only average ATT and ATC according to the probability of treatment.
%\begin{align*}
%ATE = ATT \cdot Pr(T=1) + ATC \cdot Pr(T=0)
%\end{align*}
%In the social pressure example, this would involve weighting ATT according to the number of individuals that received the mailing and weighting ATC according to the number of individuals that did not receive the mailing. This is known as the Average Treatment Effect (ATE), and can be alternatively described as the average of the ITEs for the entire population ($\frac{1}{N}\sum_{i=1}^N y_i^1- y_i^0$ where $N$ is the population size). See GH pg. 173 for an explanation with a slightly different focus. 

%\subsection{Identification of the Average Treatment Effect with Randomized Experiments}

%One reason for the disciplinary focus on average effects is that they can be estimated with randomized experiments and we only need to assume SUTVA (and even this can be relaxed in some cases). Furthermore, if we understand why randomized experiments allow us to estimate average effects, we can use this knowledge to inform our interpretations of non-randomized studies. 

%\begin{quote}
%{\it A good way to do research is to imagine the experiment that we would like to conduct, and then assess how close our data get us to that experiment.}
%\end{quote}
%\begin{center} I'm paraphrasing someone, but I can't remember who\end{center}

%To understand why randomization allows the estimation of average effects, we must utilize some mathematical notation.  
%Recall that the ATE can be written as, $\frac{1}{N}\sum_{i=1}^N (y_i^1- y_i^0)$. Because this is an average, it can be re-written in the following manner:
%\begin{align*}
%\frac{1}{N}\sum_{i=1}^N (y_i^1- y_i^0) &= \frac{1}{N}\sum_{i=1}^N y_i^1 - \frac{1}{N}\sum_{i=1}^Ny_i^0
%\end{align*}
%In the context of the social pressure example, this means that the average effect of the social pressure mailing is equal to the average outcome ($\frac{1}{N}\sum_{i=1}^N y_i^1$) that we would have observed if all individuals were sent the mailing, minus the average outcome ($\frac{1}{N}\sum_{i=1}^N y_i^0$) that we would have observed if all individuals were not sent the mailing (for the social pressure example, average outcomes are just turnout proportions).

%It is important to note that this decomposition is not possible for other summaries of the distribution of effects. For example if we wanted to know the median effect ($median(y_i^1- y_i^0)$), it cannot be broken up in this manner.
%\begin{align*}
%median(y_i^1- y_i^0) &\ne median(y_i^1) - median(y_i^0)
%\end{align*}
%Because the average effect can be decomposed in the above manner, and other summary statistics cannot, for the rest of this course, we will focus on average effects.\footnote{In fact, this lack of decomposability for summaries of the distribution of effects other than the average (e.g., the median effect, the lower and upper quartile of effects, the variance of effects, etc.) has important implications for the limits of causal knowledge.} In order to simplify notation, we will write our averages using expectation notation. That is, we will write $E[y^1]=\frac{1}{N}\sum_{i=1}^N y_i^1$ and $E[y^0]=\frac{1}{N}\sum_{i=1}^N y_i^0$. Therefore, the ATE can be written as $E[y^1] - E[y^0]$.

%If we denote the treatment variable (whether an individual received a mailing) with $T$, then the observed outcome can be written in terms of the potential outcomes and the treatment variable $Y_i = y_i^1 \cdot T_i + y_i^0 \cdot(1-T_i)$. Furthermore, randomization of treatment assignment justifies the ignorability assumption (GH pg. 183):
%\begin{align*} 
%y^0,y^1 \bot T
%\end{align*}
%In words, this means that the treatment assignment is independent of the potential outcomes. That is, knowing whether an individual received treatment (or control) does not tell us anything about their potential outcomes. This implies that the following equalities hold: 
%\begin{align}
%E[y^1] = E[Y|T=1] \label{eq:ig1} \\
%E[y^0] = E[Y|T=0] \label{eq:ig0}
%\end{align}
%In words, the first equation states that the expected (average) outcome in the population if \emph{everyone} received the treatment is equal to the expected (average) outcome in the population among those that \emph{did} receive the treatment, and the second equation states that the expected (average) outcome in the population if \emph{no one} received the treatment is equal to the expected (average) outcome in the population among those that \emph{did not} receive the treatment.

%We can now explain why randomized experiments allow the identification of average effects. For a large population, if we randomly assign half of the population to receive $T=1$, then the subpopulation with $T=1$ is just a random sample from the population, and if half of the population is large (think $\frac{N}{2} \rightarrow \infty$), we would expect the average $y^1$ from this subpopulation to equal the average $y^1$ in the overall population.\footnote{It is the consistency of the sample average that allows us to reach this conclusions.} Similarly, if we randomly assign half of the population to receive $T=0$, then the subpopulation with $T=0$ is just a random sample from the population, and if half of the population is large, we would expect the average $y^0$ in this subpopulation to equal the average $y^0$ in the overall population.\footnote{For a small population, the equalities in \eqref{eq:ig1} and \eqref{eq:ig0} will not hold exactly under the interpretation stated above. For example, suppose our population is the 50 states in the US, the outcome variable is statewide voter turnout, the treatment condition is having election day registration (EDR) available, and control is not having EDR available. In this case, even if we could randomly assign 25 of the states to treatment and 25 of the states to control, it is unlikely that the average voter turnout among the treated states \emph{for any given randomization} would equal the average voter turnout for all 50 states if all states received the treatment. The equalities in \eqref{eq:ig1} and \eqref{eq:ig0} will hold exactly across \emph{all possible randomizations}, but we would have to account for the additional variance across randomizations. This will be discussed in the next section. }

%More formally, if we randomize all units in the population to either $T=0$ or $T=1$, then we know that $T$ is independent from $y^1$ and $y^0$ (recall that independence between two random variables means that knowing the value of one does not provide information about the other). Independence between $T$ and $y^1$ means that the marginal distribution of $y^1$ ($f(y^1)$) is equal to the distribution of $y^1$ conditional on $T=t$ ($f(y^1|t)$). One consequence of this equality ($f(y^1) = f(y^1|t)$), is that Assumption \eqref{eq:ig1} holds. The explanation for Assumption \eqref{eq:ig0} is analogous.

%{\bf Summary:}  \\
%In order to define Individual Treatment Effects (ITEs), we must assume that SUTVA holds. To identify the Average Treatment Effect (ATE), we must also assume that the ignorability assumptions in \eqref{eq:ig1} and \eqref{eq:ig0} hold.

%{\bf Discussion Questions:} 
%\begin{itemize}
%\item How well does SUTVA seem to hold for Kuklinski et al. (1997), Gerber, Green, and Larimer (2008), and Fish (2002)
%\item How well does ignorability hold for Washington (2008)... 
%\end{itemize}



%\subsection{Estimation and Inference for Average Treatment Effects using Linear Regression}

%If SUTVA and the ignorability assumptions in \eqref{eq:ig1} and \eqref{eq:ig0} hold, then we can estimate ATE using linear regression. However, all of the assumptions we discussed in lectures 5, 6, and 7 still apply. We previously used the assumptions in lectures 5, 6, and 7 to make predictive (i.e., non-causal) inference. To make causal inference, we must first assume that SUTVA and ignorability hold. That allows us to interpret our predictive inferences as causal inferences. Then, we still have to go through the routine of making predictive inferences.  The following paragraphs summarize that process for a simple regression with a dichotomous treatment. 

%In order to obtain an unbiased estimate of ATE with linear regression, we must assume that $E[Y|T=t]$ can be modeled as a linear function of $t$.
%\begin{align}
%E[Y|T=t] = \beta_0 + \beta_1 \cdot t
%\end{align}
%This assumption is trivially satisfied when $T$ is binary. Note that when SUTVA, ignorability, and linearity of conditional expectation are satisfied, $\beta_1$ is the average treatment effect. 
%\begin{align*}
%ATE &= E[y^1 - y^0] \\
%&= E[y^1] - E[y^0] \\
%&= E[Y|T=1] - E[Y|T=0] \\
%&= [\beta_0 + \beta_1] - [\beta_0] \\
%&= \beta_1
%\end{align*}
%If we can obtain an unbiased estimate of $\beta_1$, we will have an unbiased estimate of ATE. Therefore, in addition to SUTVA, ignorability, and linearity of conditional expectation, if we have a representative sample from the population (or if our sample is the population), then the least squares estimator ($\widehat{\beta}_1$) provides an unbiased estimator of $\beta_1= ATE$.\footnote{If we have actually sampled the entire population (i.e., we have a census), then the population may not be very large and hence we may worry that the equalities stated in \eqref{eq:ig1} and \eqref{eq:ig0} will only hold across all possible randomizations (e.g., as in the example presented in fn. ?). However, it turns out that if the sample/population is not too small, treating the sample as if it had been randomly sampled from a very large population provides a good approximation to what we would get if we performed the correct finite-population inference (where the randomness comes from the randomized treatment assignment). For example, even if the population is the 50 states, treating those states as if they had been randomly sampled from some infinite superpopulation of states provides a good approximation of performing exact finite-population inference where we would take into account the randomness inherent in randomizing 25 of the states to receive EDR. We do not need to believe that the superpopulation of states is a real thing, it merely functions to provide an approximation.}







\end{document}


